\pdfvariable suppressoptionalinfo 611
\pdfvariable trailerid {[<C2572026083100000000000000000000><C2572026083100000000000000000000>]}
\providecommand{\CRevisionRound}{2}
\documentclass[10pt]{article}
\usepackage[a4paper,margin=17mm]{geometry}
\usepackage{amsmath,amssymb,amsthm,booktabs,microtype}
\usepackage[T1]{fontenc}
\usepackage{lmodern}
\usepackage[hidelinks,hypertexnames=false]{hyperref}
\emergencystretch=2em
\newtheorem{theorem}{Theorem}
\newtheorem{proposition}[theorem]{Proposition}
\newcommand{\C}{\mathbb C}
\newcommand{\PP}{\mathbb P^1(\C)}
\newcommand{\AM}{\mathrm{AM}}
\title{Global Quadratic Newton Dynamics through the Cayley Coordinate}
\author{HCS Research Program}
\date{31 August 2026\quad (revision \CRevisionRound)}
\begin{document}
\maketitle
\begin{abstract}
For Newton iteration of \(p_a(z)=z^2-a^2\), \(a\ne0\), we take an all-state
step: a Cayley coordinate conjugates the rational map on the Riemann sphere
exactly to squaring.  This yields both root basins, their Julia boundary,
double-exponential error identities, and a complete periodic/preperiodic
classification.  A root of unity of order \(2^e q\), \(q\) odd, has exact
tail \(e\) and eventual period \(\operatorname{ord}_q(2)\).  We also derive all
multipliers, \(\#\operatorname{Fix}(N_a^n)=2^n+1\),
\(\zeta_{\AM}(t)=((1-t)(1-2t))^{-1}\), and the invariant Cauchy law on the
Julia line.  Exact independent checks cover 16 periods, 128 root orders, and
41 hostile mutations.  The result is source-local complex dynamics, with no
arithmetic, target-determinant, or Hilbert--Pólya claim.
\end{abstract}

\section{Frozen Newton family and global coordinate}
For \(a\in\C^*\), let
\begin{equation}
 p_a(z)=z^2-a^2,\qquad
 N_a(z)=z-\frac{p_a(z)}{p_a'(z)}=\frac{z^2+a^2}{2z}
 \label{eq:newton}
\end{equation}
act on \(\PP\).  The clock is the Newton iterate \(n\ge0\).  Introduce
\begin{equation}
 C_a(z)=\frac{z-a}{z+a},\qquad
 C_a^{-1}(w)=a\frac{1+w}{1-w}.
 \label{eq:cayley}
\end{equation}
Direct cancellation gives the decisive identity
\begin{equation}
 C_a\!\left(N_a(z)\right)=C_a(z)^2,
 \qquad C_a\!\left(N_a^n(z)\right)=C_a(z)^{2^n}.
 \label{eq:conjugacy}
\end{equation}
Thus the statement below is global on the sphere; it is not a local
convergence estimate inferred from samples.

\begin{theorem}[global basin and error atlas]
The basins of \(+a\) and \(-a\) are
\begin{align}
 \mathcal B(+a)&=\{z:|C_a(z)|<1\}=\{z:\Re(z/a)>0\},\\
 \mathcal B(-a)&=\{z:|C_a(z)|>1\}=\{z:\Re(z/a)<0\}.
\end{align}
Their common boundary
\(J(N_a)=\{\Re(z/a)=0\}\cup\{\infty\}\) is the Julia set.  If
\(w=C_a(z)\), then in the plus basin
\begin{equation}
 N_a^n(z)-a=\frac{2a w^{2^n}}{1-w^{2^n}}.
 \label{eq:pluserror}
\end{equation}
If \(v=w^{-1}\), then in the minus basin
\begin{equation}
 N_a^n(z)+a=-\frac{2a v^{2^n}}{1-v^{2^n}}.
 \label{eq:minuserror}
\end{equation}
\end{theorem}

\paragraph{Proof.}
The inequalities are \(|z-a|<|z+a|\) and its reverse, which are equivalent
to the sign of \(\Re(z/a)\).  The unit circle is the Julia set of \(w^2\), so
its inverse Cayley image is the displayed generalized line.  Substituting
\(w^{2^n}\) into \(C_a^{-1}\) proves
(\ref{eq:pluserror})--(\ref{eq:minuserror}).

\section{Every periodic orbit and the dynamical zeta}
Write \(D(w)=w^2\).  The equation \(D^n(w)=w\) contains \(0\), \(\infty\),
and the \(2^n-1\) roots of \(w^{2^n-1}=1\).  Hence
\begin{equation}
 F_n=\#\operatorname{Fix}(N_a^n)=2^n+1,
 \quad P_n=\sum_{d\mid n}\mu(n/d)(2^d+1),
 \quad O_n=P_n/n.
 \label{eq:counts}
\end{equation}
Here \(P_n\) and \(O_n\) count exact-period points and primitive cycles.
The two root fixed points are superattracting.  At any Julia point of exact
period \(n\),
\((D^n)'(w)=2^n w^{2^n-1}=2^n\), so every such cycle is repelling and its
multiplier is exactly \(2^n\).  Therefore
\begin{equation}
 \zeta_{\AM}(t)=\exp\!\left(\sum_{n\ge1}\frac{F_n}{n}t^n\right)
 =\frac{1}{(1-t)(1-2t)}.
 \label{eq:zeta}
\end{equation}

\ifnum\CRevisionRound>0
\section{Exact tails and the invariant failure-line law}
The count formula does not by itself classify strictly preperiodic points.
Equation (\ref{eq:conjugacy}) supplies the missing global statement.
\begin{theorem}[complete preperiodic classification]
Apart from \(w=0,\infty\), a point is preperiodic if and only if \(w\) is a
root of unity.  If its order is
\(m=2^e q\) with \(q\) odd, then the exact tail is \(e=v_2(m)\), the landing
order is \(q\), and the eventual exact period is
\(\operatorname{ord}_q(2)\), with \(\operatorname{ord}_1(2)=1\).  It is
periodic precisely when \(e=0\).
\end{theorem}
Indeed, each squaring removes exactly one factor \(2\) from the order until
the odd order \(q\) remains.  Thereafter squaring acts by multiplication by
\(2\) in \((\mathbb Z/q\mathbb Z)^*\).  Conversely, equality between two
distinct powers \(w^{2^r}=w^{2^s}\) forces \(w\) to be a root of unity.

On the Julia set write \(w=e^{i\theta}\) and
\begin{equation}
 z=ia s,\qquad s=\cot(\theta/2).
\end{equation}
Normalized Haar measure becomes
\begin{equation}
 d\nu(s)=\frac{ds}{\pi(1+s^2)},
 \label{eq:cauchy}
\end{equation}
and angle doubling becomes the rational map
\begin{equation}
 s\longmapsto\frac{s^2-1}{2s}.
 \label{eq:boundary}
\end{equation}
Thus \(\nu\) is invariant and mixing, and the Lyapunov exponent in the
Cayley angle is \(\log2\).  In Newton language, the basin-separating line
carries an exact Cauchy statistical law rather than convergence to a root.

\paragraph{Parameter boundaries.}
Scale covariance is \(N_a(au)=aN_1(u)\), and replacing \(a\) by \(-a\)
exchanges the root labels.  At \(a=0\), however, \(N_0(z)=z/2\) for finite
\(z\): the rational degree drops from two to one.  The degree-two theorem is
not extended silently across this face.
\fi

\ifnum\CRevisionRound=1
\paragraph{Exact receipt.}
The canonical JSON records periods \(1\) through \(16\), root-of-unity orders
\(1\) through \(128\), ten exact real basin rows, and eight rational Cauchy
rows.  The producer-independent checker reconstructs all integer formulas;
SymPy separately verifies the rational conjugacy, error identities,
multipliers, boundary map, and zeta coefficients.  These finite rows are
regression probes for the analytic theorem, not its proof.
\fi

\ifnum\CRevisionRound>1
\section{Independent gates, collision control, and route boundary}
The checker passes 1317 assertions.  A separate SymPy program passes 88
identities, clean-process evidence replay is byte-identical, and all 41
repaired-hash hostile mutations are rejected.  Three revision PDFs are built
twice under a fixed epoch and checked for embedded subset fonts, clean logs,
extractable text, and visual integrity.

The ownership boundary is explicit.  C141 concerns the Hardy-space
inverse-branch Ruelle operator for \(z^2-6\), not Newton root basins.  C177
concerns Wold and Sobolev mixing for every expanding circle degree, not the
Newton sphere, root errors, or exact even-order tail.  C257 does not claim
either neighboring operator theorem; this is workspace bookkeeping, not a
literature-priority claim.

The locked scope is \texttt{NO\_BAD\_EULER\_OR\_ROOT\_NUMBER}.  Formula
(\ref{eq:zeta}) is the generic degree-two source dynamical zeta, not an
arithmetic Euler product.  There is no prime carrier, logarithmic prime clock,
arithmetic local datum, root number, automorphy statement, target divisor or
functional equation, target determinant, or Hilbert--Pólya operator.  The
strict tuple is
\texttt{(A0\_FAIL,A1\_WEAK,A2\_FAIL,A3\_FAIL,A4\_FORMAL\_HINT)}, the
overall verdict is \texttt{ROUTE\_A\_REJECTED}, and Route B is false.
\fi

\section{Conclusion}
One Möbius conjugacy closes an unusually broad dynamical ledger: global root
selection, exact quadratic error at every iterate, all periodic and
preperiodic points and multipliers, primitive counts, zeta, and the invariant
law on the failure line.  The same completeness sharpens the obstruction:
the formulas are universal for quadratic Newton maps with two simple roots
and do not single out arithmetic primes or a target spectrum.  No statement
for cubic or higher-degree Newton maps is inferred.

\begin{thebibliography}{9}
\bibitem{Cayley1879} A.~Cayley, ``Desiderata and Suggestions: No.~3. The
Newton--Fourier Imaginary Problem,'' \emph{American Journal of Mathematics}
2 (1879), 97, DOI
\href{https://doi.org/10.2307/2369201}{10.2307/2369201}.
\bibitem{ArtinMazur1965} M.~Artin and B.~Mazur, ``On Periodic Points,''
\emph{Annals of Mathematics} 81 (1965), 82--99,
\href{https://www.jstor.org/stable/1970384}{JSTOR 1970384}.
\end{thebibliography}
\end{document}
